\section{Related work}\label{sec:related}

\subsection{The recognition--production gap}\label{related:gap}

A robust body of work establishes that what a model can \emph{recognize} (rank
correctly among options, or encode in its hidden states) outruns what it can
\emph{produce} (generate correctly), and that this gap is real rather than an
evaluation artifact. Probing work
shows models often encode the correct answer in their representations while
generating an incorrect one \citep{LLMsKnowMore}; the original
inference-time-intervention work reports a large gap between a truthfulness probe's
accuracy and the model's generation accuracy \citep{ITI}; and recent evaluation work
argues that answer-matching on free generation is a stricter, more faithful test
than multiple-choice scoring \citep{AnswerMatching}. Two recent results establish
the specific, causal-and-geometric form of the gap that we take as setup: a correctness
probe reaches 0.80--0.97 AUC across eleven models, and on the subset where
output-based readouts are measured against it those readouts (P(True), token entropy,
semantic entropy) reach only 0.44--0.64, with the direction validated causally by
activation steering --- it shifts hallucination rates by $9.1\pp$ on Qwen2-7B, while
norm-matched random and orthogonal directions move them by $0.8$ and $0.3$ points and
neither reaches significance \citep{ConfidenceManifold}; and a formal
internal-vs-external knowledge measure finds an average 40\% hidden-knowledge gap
across open-weight models \citep{InsideOut}. Two scope conditions travel with the
first of these and we state them because our own setup inherits them: its probes are
teacher-forced rather than read off cold multiple-choice scoring, and its authors
report the internal advantage as regime-specific --- large on adversarial
misconceptions, closing on standard QA --- so what it establishes for us is that
correctness is linearly decodable where an output-level readout is not, not a general
ranking of internal over external estimators. We take this lossy-readout fact as \emph{established setup and
claim no priority on it} --- it is the backdrop, not our contribution. Our per-item
arbiter (the base solves under chain-of-thought $124/124$ of the items the adapter
fixes in cold scoring) confirms the gap in our channel. It also shows that between
\emph{faces} the asymmetry runs the other way --- generation succeeds precisely where
the cold discriminative readout fails --- which is the direction the generative-AI
paradox reports: generative capability is not contingent upon, and can therefore
exceed, the ability to understand the same kind of output, with models outperforming
humans at generation while falling short of them on measures of understanding
\citep{GenerativeAIParadox}. Which capacity outruns which is thus a question about the
pair of channels being compared, not a fact about models in general. Our contribution is the
\emph{converse} direction and a \emph{quantified causal attribution}: we show that a
cheap fine-tune recovers most of the cold-scoring deficit along a learned, off-axis
direction, and that the recovery does not transfer back to generation --- a property
of what the adapter does, not only of the gap's existence. Closer to our read-only
selector specifically, the gap has already been \emph{operationalized} as an
answer-selection method: an activation-derived readout selects the multiple-choice
answer and recovers an accuracy lift with no weight edit --- by attention-head-norm
voting \citep{NoVo} or a frozen correctness probe re-scoring per-option
probabilities \citep{CORAL}. We claim no priority on that capability either; we use
it only as the read leg of a read${\neq}$write dissociation (\cref{disc:selector}),
where the novelty is the contrast with an off-axis \emph{write} these works do not
study. A complementary line \emph{does} study a causal write, but in the generation
face rather than our cold-scoring channel: instruction-tuned models often fix their
answer before reasoning, and a linear probe on the pre-CoT last-token activation
decodes that answer at ${\sim}0.9$ AUC while steering along the probe direction
flips the produced answer in over 50\% of cases, exceeding orthogonal baselines
\citep{PreCoTAnswer}. This sharpens rather than blurs our claim: in that
decision/generation face an answer direction is genuinely writable, whereas in our
cold-MC scoring channel the correctness direction is decodable ($0.955$,
\cref{disc:selector}) yet the on-axis write fails ($\Wknow$ $6.4\%$,
\cref{disc:geometry}) and the off-axis write that does lift never transfers back to
generation (\cref{disc:arbiter}) --- so our read${\neq}$write dissociation is
\emph{channel-specific}, not a claim that answer directions are never writable. Its
two steering failure modes when a wrong answer is induced --- non-entailment and
confabulation, i.e.\ reasoning from a false belief --- also parallel the
reasoning-chain breakage our contrastive injection produces on gsm8k
(\cref{disc:pathologies}).

\subsection{Multiple-choice scoring of instruction-tuned models is an unreliable
readout}\label{related:mcunreliable}

That cold multiple-choice scoring is a poor proxy for an instruction-tuned model's
behavior is established, and we claim no priority on it. First-token /
log-likelihood option scores are severely misaligned with the model's actual text
answer --- mismatch rates exceeding 60\%, worst on models heavily fine-tuned on
conversational or safety data \citep{MyAnswerIsC} --- and the text answer is the
more robust quantity when the two disagree \citep{LookAtTheText}. Cold MC scoring is
further sensitive to option ordering and symbol identity
\citep{PriDe,OrderSensitivity}, is poorly calibrated without correction
\citep{CalibrateBeforeUse}, and instruction tuning polarizes the answer distribution
\citep{folktexts}. One recent result is especially close to our framing: the chat
template itself inflates a model's confidence in ``its own'' answer in a way
removable by a parameter-free reframing \citep{OwnershipBias}. A second is close in
mechanism but sits on the opposite population, and we flag the mismatch rather than
borrow the conclusion: on \emph{intermediate pretraining checkpoints} --- small,
\emph{not} instruction-tuned models --- emitting the option symbol is a skill
separable from knowing the answer, and fine-tuning on multiple-choice format first
disentangles the evaluation of knowledge from the emerging ability to follow the
format \citep{FormatFirst}. Our model is instruction-tuned, i.e.\ the population that
work treats as already past the difficulty, so it does not license the claim that our
base cannot emit the symbol --- and indeed ours can, which is why the one-forward
letter-emission baseline is so strong (\cref{disc:selector}). What transfers is the
weaker and still useful point: symbol emission and knowledge are separable enough
that a protocol can score one while intending the other. Selection bias is moreover introduced by supervised
fine-tuning itself, tracing to weak multiple-choice symbol binding --- the A/B/C/D
symbol--content association --- which strengthening during SFT both reduces and turns
into an accuracy gain \citep{SymbolBinding}; this bears directly on our matched-SFT
control and on the trivial letter-emission baseline we find outscores the adapter
(\cref{disc:selector}). We build directly on this literature --- the
protocol (cold MC under a chat template) is a lossy readout, which we treat as
\emph{established setup rather than as our thesis} --- and our contribution is past
the diagnosis: we quantify how much of the deficit a generic fine-tune recovers
(${\sim}62\%$ on ARC, a matched-control attribution), locate that recovery
geometrically (off-axis of degree, established causally), and show it does not
transfer. Critically, and unlike a ``the lift is a formatting artifact'' reading,
the recovered lift is \emph{real} signal: it survives a same-scaffold PMI
decomposition (ARC retains 78\%), so the readout is genuinely leaving knowledge on
the table, not merely mis-normalizing it.

\subsection{The mirror in compression: the benchmark illusion}
\label{related:benchillusion}

The closest single neighbor in the same channel (log-likelihood scoring vs greedy
generation) comes from model compression. Pruned models can pass multiple-choice
evaluation while failing to answer the same question in open generation; the correct
answer is usually not erased but \emph{demoted} in rank, reappearing under beam
search, sampling, or one in-context example \citep{BenchmarkIllusion}. That is our
phenomenon's mirror: there, MC scoring \emph{overstates} usable capability (the
answer is demoted but MC-recoverable); here, cold MC scoring \emph{understates} it
(the answer is CoT-accessible but cold-misranked) and a fine-tune \emph{promotes} it
in the MC channel without helping generation. The shared lesson --- multiple-choice
scores are not generation capability --- is established; the constructive,
opposite-signed instance, and its off-axis mechanism, are ours.

\subsection{The knowledge--prediction gap and KAPPA (our foil, not our
competitor)}\label{related:kappa}

Most directly adjacent is work on the knowledge--prediction gap in multiple-choice
questions: models fail MCQs whose answers they encode internally, and KAPPA
identifies distinct knowledge and prediction subspaces in the residual stream and
aligns them with a closed-form, inference-time intervention \citep{KAPPA}. We
position KAPPA as a foil rather than a competitor, and explicitly \emph{not} as a
head-to-head baseline we beat. KAPPA operates where the correct option is present
and decodable in the residual (options-present MCQ); our effect lives in cold
options-absent text-continuation scoring, where we find no evidence that the
precondition holds --- a
faithful $k$-class knowledge probe in our channel decodes the gold option at
$0.317$ against chance $0.25$, while the same probe for the greedy option
decodes at $0.542$. We state this as a \emph{limit on power rather than a
demonstration of absence}, because the design cannot separate the two: the gold
figure rests on a $120$-item dev split, whose Wilson $95\%$ interval
$[0.240, 0.404]$ contains chance \emph{and} contains improvements larger than
$10$\,pp, so an equivalence test at that margin does not pass; and the point itself
is a maximum over a $40$-cell (layer $\times$ regularization) grid, which biases it
upward exactly where it is being read as a null. What we can say is that KAPPA's
mechanism is not \emph{detectably} instantiated in our channel at this resolution;
the relationship is a channel mismatch, and we make no claim of out-performing it.
One apparent tension deserves a sentence here, because a reader who reaches
\cref{disc:selector} will otherwise trip on it: our read-only selector decodes the
gold at \textbf{0.943}, far above this $0.317$. The two probes read \emph{different
channels}. The $0.317$ is a $k$-class probe on the \emph{prompt-final residual} ---
a single vector with the options in context, which is precisely the channel KAPPA
operates on --- whereas the $0.943$ selector reads the \emph{per-head activations of
each option continuation} (four separate forwards, arg-max over the per-option
score). KAPPA's precondition fails in our channel, and the selector does not restore
it; it succeeds by reading somewhere else. We further reconcile KAPPA's reported
free-form \emph{generalization} with our no-transfer result
(\cref{disc:pathologies}), since the two superficially conflict: KAPPA's own
free-form evaluation carries its gain almost entirely on a single
answer-prior-shaped task (TruthfulQA $+2.5$) while it \emph{regresses} on reasoning
generation (GSM8K $91.6{\rightarrow}90.7$) and is flat on BBQ-Age ($+0.2$) --- so
KAPPA, like our lift, does not transfer to reasoning generation, which corroborates
rather than contradicts our finding.

\subsection{Answer priors and surface-form competition}\label{related:answerpriors}

Our task-heterogeneity finding --- that the TruthfulQA half of the lift is largely
an answer-prior / surface-form effect --- rests on a well-developed line.
Surface-form competition shows that highest-probability scoring is biased by how
answers are phrased, motivating PMI-style domain-conditional normalization
\citep{SurfaceFormCompetition}; choices-only probing shows models answer MCQs above
chance with the question removed \citep{ArtifactsOrAbduction,ChoicesOnlyCheater};
and choice-sensitivity admits parameter-free fixes \citep{NPSQ}. We use the same
PMI/decision-conditional and choices-only controls to separate the two tasks: ARC is
question-conditional and carries the mechanism, whereas TruthfulQA is largely
answer-prior, which is why we report it held-out as evidence the lift is real and
never as evidence for the off-axis mechanism.

\subsection{The capability tax of interventions, and the specificity question}
\label{related:tax}

Adjacent work shows that adapting a model along one axis taxes others: knowledge
injection under PEFT erodes reasoning, with the reasoning-critical information
distributed across the singular spectrum rather than localized \citep{PALoRA}, and
supervised fine-tuning degrades pretrained capabilities because its targets diverge
from the model's autoregressive distribution \citep{MixSD}. Inference-time steering
can be applied selectively, choosing when and how much to intervene \citep{MERA}; and steering vectors that cleanly control a target property in
multiple-choice or toy settings degrade in free-form generation, inducing degenerate
repetition and factual hallucination at the strengths control requires
\citep{BeyondMultipleChoiceSteering} --- a control--quality trade-off that mirrors the
creative degeneration our contrastive write produces (\cref{disc:pathologies}).
These motivate our deployment-pathology results, but they also raise the sharpest
reviewer objection to an off-axis claim: that \emph{any} norm-matched perturbation
works, so an off-axis direction is not identifiable \citep{NonIdentifiability}, and
even random directions can move behavior \citep{RogueScalpel}. We answer this
directly, and with the same pair of controls that work runs on its own steering
result --- the learned direction against a norm-matched random one and against an
orthogonal one, neither of which reaches significance there
\citep{ConfidenceManifold}: a norm-matched random offset
recovers ${\approx}0\%$ of our lift (so the direction is learned, not arbitrary),
the strongest on-axis reference recovers $6.4\%$, and the trained offset has an
optimal magnitude (an inverted-U strength sweep), which a generic perturbation would
not. We also note that reported specificity failures of activation directions are
cross-model; within-model --- our setting --- directions remain specific
\citep{SyedActivationDirections}, so that result does not subsume our control. A
related limit on the \emph{write} side sharpens what an off-axis control claim must
contend with: a dual-stance evaluation of a sycophancy-reduction direction finds
that although sycophantic and factual agreement occupy geometrically \emph{distinct}
subspaces, a single steering direction projects equally onto both and cannot
differentially target either --- a case where representations readable from
activations are not separately writable through them \citep{DualStanceSycophancy}.
That is the writability gap our read${\neq}$write dissociation
(\cref{disc:selector,disc:geometry}) makes explicit from the opposite side: there a
writable direction fails to \emph{separate} two readable targets, whereas here the
on-axis correctness direction is readable yet fails to \emph{drive} the lift
($\Wknow$ $6.4\%$) while an off-axis one does. A distinct mechanistic hypothesis for
any logit-level intervention is temperature regulation --- confidence-regulation /
entropy neurons write into the unembedding null space to rescale the logit
temperature without re-ordering the argmax \citep{ConfidenceRegulationNeurons} ---
which we test and rule out for our offset in \cref{disc:negatives}.

\subsection{PEFT, representation editing, and the fine-tuning baseline}
\label{related:peft}

Our adapter is from the localized-fine-tuning / representation-intervention family:
LoFiT learns additive per-head offset vectors on a sparse set of attention heads
\citep{LoFiT}, a descendant of inference-time intervention \citep{ITI}, and related
to representation fine-tuning \citep{ReFT}; concept- and feature-level steering
benchmarks find that fine-tuning and prompting often outperform sparse-feature
methods \citep{AxBench}. A recent principled account establishes a first-order
equivalence between activation-space interventions and weight-space updates and finds
the two play distinct, complementary functional roles
\citep{WeightUpdatesActivationShifts}; this is why we compare methods by their
\emph{net effect on activations} rather than by the weight-space offset, which we show
is a misleading proxy (\cref{disc:geometry}). The decisive baseline our inducibility
result requires is
ordinary supervised fine-tuning: prior domain-adaptation work establishes SFT as a
strong, cost-effective default for multiple-choice accuracy while degrading
open-ended generation \citep{FrenchQASFT}. Our matched Align-LoRA control
instantiates exactly this baseline, and our finding --- that it recovers
${\sim}62\%$ of the contrastive lift --- places most of the effect on fine-tuning in
general rather than on the contrastive objective.
